% GENERATED FILE. Edit canonical lessons, then run npm run generate:books.
% canonical-source-hash: fnv1a64:c4d090a7c81fbd55
% canonical-lessons: LA-C325-bulla, LA-C325-crepundia, LA-C325-praenomen, LA-C325-cognomen, LA-C325-tutor

\chapter{Nōmen et bulla}
\label{ch:la-nomen-et-bulla}
\hlchaptermodality{325}

\begin{chapteropening}
\textbf{By the end of this chapter:} \emph{I can talk about a Roman child growing up: the amulet, the toys, the first name, the extra family name and the guardian.}

Complete the last lesson of chapter 325: I can talk about a Roman child growing up: the amulet, the toys, the first name, the extra family name and the guardian.
\end{chapteropening}

\section[bulla, bullae]{bulla, bullae — an amulet (worn by a child); a bubble}

\label{lesson:LA-C325-bulla}



\begin{quote}
Before the new one: say the Latin for childhood, then the Latin for to bring up.
\end{quote}

\subsection*{You'll want to know: bulla, bullae}
\textbf{bulla, bullae} — \textquotedblleft{}an amulet (worn by a child); a bubble\textquotedblright{}.

\begin{quote}
Puer bullam auream in collō gerēbat.
\end{quote}

\begin{quote}
\emph{The boy wore a golden amulet round his neck.}
\end{quote}

Feminine. First a bubble, then the round locket a freeborn boy wore until he came of age. English \emph{bullet} and the papal \emph{bull} come from it.

\begin{cousinweb}
\textbf{bulla} first meant \textquotedblleft{}a bubble\textquotedblright{}, anything swelling into a round shape (Lewis \& Short); where the word comes from is not settled. English \emph{bull}, a papal letter with a round seal, comes from it, and \emph{bullet} through French.
\end{cousinweb}

\begin{grammarlens}[title={What you've built}]
One more word for telling, explaining and describing at B1.
\end{grammarlens}

\subsection*{Your turn}
\begin{itemize}
\raggedright
  \item \emph{Say it:} \emph{bulla, bullae}
  \item \emph{Say it:} \emph{bulla, bullae}, once more
  \item \emph{Say it:} a sentence of your own with \emph{bulla, bullae}
  \item \emph{Recall:} say the Latin for childhood, then the Latin for to bring up, then say \emph{bulla, bullae} again
\end{itemize}

\begin{tcolorbox}[breakable,colback=teal!4,colframe=teal!35!black,title={Before you move on}]
What does bulla, bullae mean? (\textquotedblleft{}An amulet (worn by a child); a bubble\textquotedblright{}.) Say it once more.
\end{tcolorbox}
\section[crepundia, crepundiōrum]{crepundia, crepundiōrum — a rattle; childhood toys}

\label{lesson:LA-C325-crepundia}



\begin{quote}
Before the new one: say the Latin for to bring up, then the Latin for an amulet.
\end{quote}

\subsection*{You'll want to know: crepundia, crepundiōrum}
\textbf{crepundia, crepundiōrum} — \textquotedblleft{}a rattle; childhood toys\textquotedblright{}.

\begin{quote}
Nūtrīx īnfantī crepundia dedit.
\end{quote}

\begin{quote}
\emph{The nurse gave the baby a rattle.}
\end{quote}

Neuter, plural only. In comedies a lost child is recognised by the little toys kept in a box.

\begin{cousinweb}
\textbf{crepundia} is built on \emph{crepō}, \textquotedblleft{}to rattle, to creak\textquotedblright{} (Lewis \& Short): things that rattle.
\end{cousinweb}

\begin{grammarlens}[title={What you've built}]
One more word for telling, explaining and describing at B1.
\end{grammarlens}

\subsection*{Your turn}
\begin{itemize}
\raggedright
  \item \emph{Say it:} \emph{crepundia, crepundiōrum}
  \item \emph{Say it:} \emph{crepundia, crepundiōrum}, once more
  \item \emph{Say it:} a sentence of your own with \emph{crepundia, crepundiōrum}
  \item \emph{Recall:} say the Latin for to bring up, then the Latin for an amulet, then say \emph{crepundia, crepundiōrum} again
\end{itemize}

\begin{tcolorbox}[breakable,colback=teal!4,colframe=teal!35!black,title={Before you move on}]
What does crepundia, crepundiōrum mean? (\textquotedblleft{}A rattle; childhood toys\textquotedblright{}.) Say it once more.
\end{tcolorbox}
\section[praenōmen, praenōminis]{praenōmen, praenōminis — a first name}

\label{lesson:LA-C325-praenomen}



\begin{quote}
Before the new one: say the Latin for an amulet, then the Latin for a rattle.
\end{quote}

\subsection*{You'll want to know: praenōmen, praenōminis}
\textbf{praenōmen, praenōminis} — \textquotedblleft{}a first name\textquotedblright{}.

\begin{quote}
Pater fīliō praenōmen avī dedit.
\end{quote}

\begin{quote}
\emph{The father gave his son his grandfather's first name.}
\end{quote}

Neuter, 3rd declension: the name before the family name. Romans had few first names and wrote them short.

\begin{cousinweb}
\textbf{praenōmen} joins \emph{prae}, \textquotedblleft{}before\textquotedblright{}, and \emph{nōmen}, \textquotedblleft{}name\textquotedblright{} (Lewis \& Short): the name in front of the family name.
\end{cousinweb}

\begin{grammarlens}[title={What you've built}]
One more word for telling, explaining and describing at B1.
\end{grammarlens}

\subsection*{Your turn}
\begin{itemize}
\raggedright
  \item \emph{Say it:} \emph{praenōmen, praenōminis}
  \item \emph{Say it:} \emph{praenōmen, praenōminis}, once more
  \item \emph{Say it:} a sentence of your own with \emph{praenōmen, praenōminis}
  \item \emph{Recall:} say the Latin for an amulet, then the Latin for a rattle, then say \emph{praenōmen, praenōminis} again
\end{itemize}

\begin{tcolorbox}[breakable,colback=teal!4,colframe=teal!35!black,title={Before you move on}]
What does praenōmen, praenōminis mean? (\textquotedblleft{}A first name\textquotedblright{}.) Say it once more.
\end{tcolorbox}
\section[cognōmen, cognōminis]{cognōmen, cognōminis — an extra family name, a surname}

\label{lesson:LA-C325-cognomen}



\begin{quote}
Before the new one: say the Latin for a rattle, then the Latin for a first name.
\end{quote}

\subsection*{You'll want to know: cognōmen, cognōminis}
\textbf{cognōmen, cognōminis} — \textquotedblleft{}an extra family name, a surname\textquotedblright{}.

\begin{quote}
Ob magnam victōriam cognōmen novum accēpit.
\end{quote}

\begin{quote}
\emph{Because of a great victory he received a new surname.}
\end{quote}

Neuter, 3rd declension. The third name, which often told something about a person or a branch of the family.

\begin{cousinweb}
\textbf{cognōmen} joins \emph{co-}, \textquotedblleft{}together with\textquotedblright{}, and \emph{nōmen}, \textquotedblleft{}name\textquotedblright{} (EtymDB): the name added to the family name.
\end{cousinweb}

\begin{grammarlens}[title={What you've built}]
One more word for telling, explaining and describing at B1.
\end{grammarlens}

\subsection*{Your turn}
\begin{itemize}
\raggedright
  \item \emph{Say it:} \emph{cognōmen, cognōminis}
  \item \emph{Say it:} \emph{cognōmen, cognōminis}, once more
  \item \emph{Say it:} a sentence of your own with \emph{cognōmen, cognōminis}
  \item \emph{Recall:} say the Latin for a rattle, then the Latin for a first name, then say \emph{cognōmen, cognōminis} again
\end{itemize}

\begin{tcolorbox}[breakable,colback=teal!4,colframe=teal!35!black,title={Before you move on}]
What does cognōmen, cognōminis mean? (\textquotedblleft{}An extra family name, a surname\textquotedblright{}.) Say it once more.
\end{tcolorbox}
\section[tūtor, tūtōris]{tūtor, tūtōris — a guardian, a protector}

\label{lesson:LA-C325-tutor}



\begin{quote}
Before the new one: say the Latin for a first name, then the Latin for an extra family name.
\end{quote}

\subsection*{You'll want to know: tūtor, tūtōris}
\textbf{tūtor, tūtōris} — \textquotedblleft{}a guardian, a protector\textquotedblright{}.

\begin{quote}
Avunculus tūtor puerōrum factus est.
\end{quote}

\begin{quote}
\emph{Their uncle became the boys' guardian.}
\end{quote}

Masculine, 3rd declension. A child with no father had a \emph{tūtor} to keep his property safe. English \emph{tutor}.

\begin{cousinweb}
\textbf{tūtor} is built on \emph{tueor}, \textquotedblleft{}to watch over, to protect\textquotedblright{} (Lewis \& Short): a protector. English \emph{tutor} comes from it.
\end{cousinweb}

\begin{grammarlens}[title={What you've built}]
One more word for telling, explaining and describing at B1.
\end{grammarlens}

\subsection*{Your turn}
\begin{itemize}
\raggedright
  \item \emph{Say it:} \emph{tūtor, tūtōris}
  \item \emph{Say it:} \emph{tūtor, tūtōris}, once more
  \item \emph{Say it:} a sentence of your own with \emph{tūtor, tūtōris}
  \item \emph{Recall:} say the Latin for a first name, then the Latin for an extra family name, then say \emph{tūtor, tūtōris} again
\end{itemize}

\begin{tcolorbox}[breakable,colback=teal!4,colframe=teal!35!black,title={Before you move on}]
What does tūtor, tūtōris mean? (\textquotedblleft{}A guardian, a protector\textquotedblright{}.) Say it once more.
\end{tcolorbox}
